\documentclass[11pt]{article}

\usepackage[T1]{fontenc}
\usepackage[utf8]{inputenc}
\usepackage[margin=1in]{geometry}
\usepackage{graphicx}
\usepackage{booktabs}
\usepackage{amsmath}
\usepackage{amssymb}
\usepackage{xcolor}
\usepackage{listings}
\usepackage{caption}
\usepackage[colorlinks=true,allcolors=blue!55!black]{hyperref}

\graphicspath{{./}{figures/}}

% ---------------------------------------------------------------------------
% Figure macro: tries ./X.pdf, then figures/X.pdf, then draws a labelled box.
% Place figures in the same folder as this .tex file, or in figures/.
% ---------------------------------------------------------------------------
\newcommand{\figorbox}[2][\linewidth]{%
  \IfFileExists{#2.pdf}{%
    \includegraphics[width=#1,keepaspectratio]{#2.pdf}%
  }{%
    \IfFileExists{figures/#2.pdf}{%
      \includegraphics[width=#1,keepaspectratio]{figures/#2.pdf}%
    }{%
      \setlength{\fboxsep}{0pt}%
      \fbox{\parbox[c][0.20\textheight][c]{\dimexpr#1-2\fboxrule\relax}{%
        \centering\ttfamily\small\detokenize{#2.pdf}\\[6pt]%
        \normalfont\itshape\small copy from reports4/ to compile%
      }}%
    }%
  }%
}

% ---------------------------------------------------------------------------
% Python listing style
% ---------------------------------------------------------------------------
\lstdefinestyle{pycode}{
  language=Python,
  basicstyle=\small\ttfamily,
  keywordstyle=\bfseries\color{blue!70!black},
  commentstyle=\itshape\color{gray!80!black},
  stringstyle=\color{green!45!black},
  showstringspaces=false,
  frame=single,
  rulecolor=\color{black!18},
  backgroundcolor=\color{black!3},
  numbers=left,
  numberstyle=\tiny\color{gray},
  numbersep=7pt,
  xleftmargin=14pt,
  breaklines=true,
  captionpos=b,
  aboveskip=6pt,
  belowskip=4pt,
}

\captionsetup{font=small,labelfont=bf}
\setlength{\parskip}{3pt}
\setlength{\emergencystretch}{3em}

\title{\textbf{Designing Guide RNAs for an Entire Transposable Element Family}\\[4pt]
{\large An Algorithmic Approach to Coverage-Maximising and\\
Expression-Weighted Multi-Locus Targeting}}
\author{Anmol Dash}
\date{\today}

% ===========================================================================
\begin{document}
\maketitle

\begin{abstract}
\noindent
Silencing or activating a transposable element (TE) family with CRISPR requires choosing a
20\,bp guide RNA (gRNA) spacer that is both (i) adjacent to a protospacer adjacent motif
(PAM) in as many family members as possible, and (ii) concentrated in the loci contributing
most to transcriptional output.  These two objectives---\emph{sequence coverage maximisation}
(SCM) and \emph{expression-weighted coverage maximisation} (EWCM)---diverge whenever a
highly active subset of loci is sequence-diverged from the family majority.  We describe a
three-step algorithm: enumerate all PAM-adjacent 20-mer candidates from all locus sequences,
score them in $O(1)$ using an inverted index, and select a compact multi-guide set via
greedy approximation to the NP-hard set-cover problem.  Applied to 2,573 MT2\_Mm loci
(the LTR component of MERVL, active in the early mouse embryo) with per-locus SQuIRE
expression across five developmental stages, a single guide covers 85.7\% of loci and
91.1\% of total family expression; five guides together reach 98.8\% of loci and 99.7\%
of expression.  The algorithm, code, and figures are documented here with sufficient detail
to reproduce or adapt the approach for any TE family.
\end{abstract}

% ===========================================================================
\section{Background: Why Designing a gRNA for a TE Family Is Hard}

CRISPR--Cas9 works as follows.  A short RNA molecule---the guide RNA, or gRNA---base-pairs
with its complementary DNA sequence in the genome.  The Cas9 protein rides along and, upon
finding that match, cuts the DNA.  For silencing, a catalytically dead variant (dCas9) fused
to a repressor (KRAB) is used instead.  The only hard constraint is that the 20\,bp target
sequence must be immediately followed by a short motif called a \textbf{protospacer adjacent
motif (PAM)}: for the most widely used protein, SpCas9, the PAM is
\textbf{5$'$-NGG-3$'$} (two G's preceded by any nucleotide).  Without a PAM, Cas9 does not
bind.

For a single gene with one target, the design problem is trivial: pick any 20-mer that is
flanked by NGG.  For a TE \emph{family} the problem is combinatorial.  Consider MT2\_Mm, the
LTR element of the murine endogenous retrovirus MERVL.  There are \textbf{2,573 copies} of
this element scattered across the mouse genome, each slightly mutated relative to the
ancestral sequence---some by only a few bases, others by dozens.  A gRNA that matches one
copy perfectly may miss hundreds of others.  The question is: which 20-mer, when checked
against all 2,573 loci, appears in the largest number of them?  And can we do even better
by accounting for \emph{how much} each locus is expressed?

% ===========================================================================
\section{The Algorithm, Step by Step}

\subsection{Step 1 --- Enumerating Every Valid Target Site}

The first task is to find every genomic position in every locus where Cas9 \emph{could}
bind---i.e.\ every 20-mer that is immediately followed by NGG.  Because the genome is
double-stranded, we must check both the forward and reverse-complement strand of each locus
sequence.

\begin{lstlisting}[style=pycode,
  caption={PAM-site scanning.  For each locus sequence we slide a window of length
    $\ell + 3$ (spacer + PAM), check whether the last two bases are GG, and if so
    record the 20-mer spacer.  The reverse complement is checked identically.}]
def extract_pam_candidates(sequence, grna_len=20):
    """Return the set of unique PAM-adjacent spacers in this locus (both strands)."""
    candidates = set()
    seq = sequence.upper()
    rc  = reverse_complement(seq)     # check the other strand too

    for s in (seq, rc):
        for i in range(len(s) - grna_len - 2):
            pam    = s[i + grna_len : i + grna_len + 3]
            spacer = s[i : i + grna_len]
            # NGG: middle and last base must be G; first (N) can be anything
            if pam[1] == "G" and pam[2] == "G" and "N" not in spacer:
                candidates.add(spacer)
    return candidates
\end{lstlisting}

Running this over all 2,573 MT2\_Mm sequences finds \textbf{121,256 PAM-site positions}.
After deduplication---many copies share the same 20-mer sequence even though they sit at
different chromosomal coordinates---we are left with \textbf{6,034 unique candidate spacers}.
This roughly 20-fold compression is what makes the problem tractable.

\subsection{Step 2 --- Scoring Candidates with an Inverted Index}

A brute-force approach would, for each of the 6,034 candidates, scan all 2,573 loci to
count how many contain it.  That is $6{,}034 \times 2{,}573 \approx 15.5\,\mathrm{M}$
comparisons.  We use a faster data structure: an \textbf{inverted index}, the same idea
that makes internet search fast (a word-to-page map rather than a page-to-word map).

\begin{lstlisting}[style=pycode,
  caption={Building the inverted index.  After this runs once, any coverage query
    is an O(1) dictionary look-up.}]
from collections import defaultdict

index = defaultdict(set)   # spacer  ->  {locus_0, locus_7, locus_42, ...}

for locus_idx, row in enumerate(df.itertuples()):
    for spacer in extract_pam_candidates(row.Seq):
        index[spacer].add(locus_idx)   # record which loci contain this spacer
\end{lstlisting}

Once built, two coverage metrics follow directly:

\begin{align}
  \text{Sequence coverage}(g) &\;=\;
    \frac{\bigl|\,\mathrm{index}[g]\,\bigr|}{N}
  \label{eq:scm} \\[4pt]
  \text{Expression coverage}(g) &\;=\;
    \frac{\displaystyle\sum_{i\,\in\,\mathrm{index}[g]} e_i}
         {\displaystyle\sum_{i=1}^{N} e_i}
  \label{eq:ewcm}
\end{align}

where $N = 2{,}573$ loci and $e_i$ is the total SQuIRE count for locus $i$ summed across
all five developmental stages (pronucleus, 2-cell, 4-cell, 8-cell, morula).

\paragraph{Mismatch tolerance.}
Real Cas9 tolerates up to $\sim$3 mismatches, especially outside a 12-bp ``seed region''
immediately adjacent to the PAM where fidelity is critical.  To compute mismatch-tolerant
coverage for a guide at allowance $d$, we generate all single- or double-base variants of
the spacer that fall \emph{outside} the seed region (60 variants for $d{=}1$;
$\binom{20-12}{2}\times 9 = 252$ for $d{=}2$ in the non-seed positions), look each one up
in the index, and take the union of matched loci.  This is far cheaper than pairwise
sequence alignment across thousands of loci.

\subsection{Step 3 --- Selecting a Minimal Multi-Guide Set}

A single guide cannot reach every locus.  The best single guide covers 85.7\% of MT2\_Mm
loci; to push beyond that we need additional guides.  Finding the \emph{minimum} set of
guides that covers all (or a target fraction of) loci is equivalent to the classical
\emph{set cover} problem, which is NP-hard.  However, the greedy algorithm---at each step
choose the guide that covers the most \emph{not-yet-covered} expression---is guaranteed to
achieve $(1 - 1/e) \approx 63\%$ of the optimal coverage with $k$ guides.

\begin{lstlisting}[style=pycode,
  caption={Greedy $k$-guide selection.  At each round we pick the guide with the
    highest marginal gain over what is already covered, then zero out those loci so
    they are not double-counted.}]
def greedy_guide_set(index, expr_weights, k=5):
    uncovered = expr_weights.copy()   # copy so we can zero out covered loci
    selected  = []

    for _ in range(k):
        # find the guide that captures the most uncovered expression
        best = max(index, key=lambda g: sum(uncovered[i] for i in index[g]))
        marginal_expr = sum(uncovered[i] for i in index[best])
        new_loci      = len([i for i in index[best] if uncovered[i] > 0])
        selected.append((best, new_loci, marginal_expr))

        # mark those loci as covered
        for locus_idx in index[best]:
            uncovered[locus_idx] = 0.0

    return selected
\end{lstlisting}

In practice the curve saturates after the first two guides: guide \#1 does most of the work
and each subsequent guide adds diminishing marginal returns (Table~\ref{tab:greedy}).

% ===========================================================================
\section{Results: MT2\_Mm in the Early Mouse Embryo}

\subsection{The data}

MT2\_Mm is the long-terminal-repeat (LTR) component of MERVL, a murine endogenous
retrovirus known to be transcriptionally active during zygotic genome activation.  We used
2,573 loci annotated by RepeatMasker in \texttt{mm10}, each with a DNA sequence and
per-locus SQuIRE read counts across five stages of early mouse development:
pronucleus, 2-cell, 4-cell, 8-cell, and morula.

\subsection{Greedy guide set}

\begin{table}[h]
  \centering
  \small
  \setlength{\tabcolsep}{7pt}
  \begin{tabular}{@{}r l r r r@{}}
    \toprule
    \textbf{\#} & \textbf{Spacer (5$'$\,$\rightarrow$\,3$'$)} &
      \textbf{New loci} & \textbf{Cum.\ loci} & \textbf{Cum.\ expr} \\
    \midrule
    1 & \texttt{TGGAGACAATATAAGGACAT} & 2,204 & 85.7\% & 91.1\% \\
    2 & \texttt{ACTGGTGAGCCTTCCAATTC} &   218 & 94.1\% & 96.9\% \\
    3 & \texttt{CTGAGTAACTGCTAGATCCT} &    85 & 97.4\% & 99.2\% \\
    4 & \texttt{CACCAGTGACCCTTATCTGG} &    14 & 98.0\% & 99.5\% \\
    5 & \texttt{CTCTAGAGTCACAGAACTTA} &    22 & 98.8\% & 99.7\% \\
    \bottomrule
  \end{tabular}
  \caption{Greedy 5-guide set for MT2\_Mm (SpCas9, NGG PAM, expression-weighted).
    \emph{New loci}: loci first reached by this guide.
    \emph{Cum.\ loci / expr}: cumulative fraction of all 2,573 loci and of total SQuIRE
    expression covered by guides 1--$k$.}
  \label{tab:greedy}
\end{table}

Guide \#1 alone covers 85.7\% of loci and 91.1\% of expression---a single synthesised
oligonucleotide reaches the overwhelming majority of the family.  Adding guide \#2
pushes both figures past 94\%, and a 3-guide pool reaches $>$97\%.  Guides \#4--5
add very little: \#4 adds only 14 new loci and \#5 adds 22, suggesting those loci
carry unique 20-mer sequences not shared with the rest of the family.

\subsection{Coverage figures}

\begin{figure}[h]
  \centering
  \figorbox[\linewidth]{fig_grna_coverage}
  \caption{Guide RNA design summary for MT2\_Mm (6,034 SpCas9 NGG candidates from
    2,573 loci).
    \textbf{(a)} Distribution of exact sequence coverage across all 6,034 candidates.
    \textbf{(b)} Sequence coverage rank vs.\ expression coverage rank for every candidate;
    $r = 0.998$ (Pearson).
    \textbf{(c)} Mismatch-tolerance curve for the top EWCM guide: coverage rises as
    mismatches are allowed outside the 12-bp seed region.
    \textbf{(d)} Heatmap of the top 5 guides against HDBSCAN sequence clusters,
    showing which guides are uniformly effective and which are cluster-specific.}
  \label{fig:coverage}
\end{figure}

\begin{figure}[h]
  \centering
  \figorbox[0.82\linewidth]{fig_grna_positional}
  \caption{Density of NGG PAM sites as a function of normalised position within
    MT2\_Mm sequences (0 = 5$'$ end, 1 = 3$'$ end), shown separately for each strand.
    The approximately uniform distribution means usable target sites are spread
    throughout the element with no structural bottleneck.}
  \label{fig:positional}
\end{figure}

\begin{figure}[h]
  \centering
  \figorbox[0.88\linewidth]{fig_grna_candidates}
  \caption{Top gRNA candidates ranked by expression-weighted coverage.  The first
    guide dominates; subsequent guides provide incremental marginal gain consistent
    with the greedy saturation curve.}
  \label{fig:candidates}
\end{figure}

% ===========================================================================
\section{What the Numbers Mean Biologically}

\paragraph{85.7\% sequence coverage.}
The top guide---\texttt{TGGAGACAATATAAGGACAT}---is present as an exact 20-mer adjacent to
NGG in 2,204 of the 2,573 loci.  The remaining $\sim$14\% of copies have mutated that
stretch enough that the guide no longer matches exactly.  Allowing a single mismatch
outside the seed region (Figure~\ref{fig:coverage}c) recovers additional loci; the exact
number depends on how far those diverged copies have drifted.

\paragraph{Why SCM and EWCM agree here ($r = 0.998$).}
The Pearson correlation between sequence coverage ranks and expression coverage ranks is
very close to 1 for MT2\_Mm.  This means that the loci covered by the top guide are not
a biased sample of the expression distribution---highly expressed loci are spread throughout
the sequence-coverage landscape rather than concentrated in a diverged sub-cluster.  This
is not always the case.  For an older family in which an active sub-family has drifted
substantially from the rest, the SCM and EWCM optima can diverge: the guide covering the
most \emph{loci} may be different from the one covering the most \emph{transcription}.
The correlation coefficient is therefore a useful diagnostic: $r < 0.7$ is a signal to
prefer EWCM if the goal is to suppress transcriptional output rather than maximise
genomic targeting breadth.

\paragraph{Uniform PAM density (Figure~\ref{fig:positional}).}
PAM sites are approximately uniformly distributed along the MT2\_Mm consensus.  In some
retroviral LTR elements, PAM sites cluster near the LTR boundaries or the primer-binding
site (PBS); such clustering would constrain guide choice to specific structural domains.
The uniform density here means the algorithm has access to target sites throughout the
element and can freely choose the spacer with the best coverage properties rather than
the best \emph{available} one.

\paragraph{Practical guide for a CRISPRi experiment.}
\begin{itemize}
  \item If the goal is maximal silencing with a single reagent, order guide \#1
    (\texttt{TGGAGACAATATAAGGACAT}).  GC content = 40\%, no homopolymer runs,
    no 5$'$-T: all standard quality checks pass.
  \item If near-complete silencing is required, a 2- or 3-guide pool is sufficient;
    four or five guides add negligible additional coverage.
  \item The full ranked table (\texttt{grna\_candidates.csv}) lists all 6,034 candidates
    with sequence coverage, expression coverage, and on-target quality score, so any
    subset can be filtered by experimental criteria (e.g.\ exclude guides with $>$80\%
    seed-region GC).
\end{itemize}

\paragraph{Extending to other Cas proteins.}
The algorithm is PAM-agnostic.  Changing the Cas protein flag to \texttt{SaCas9} switches
the PAM to NNGRRT (6\,bp); \texttt{Cas12a} uses a 5$'$-TTTV PAM and cuts in the opposite
orientation.  SpRY, a near-PAMless SpCas9 variant, uses NRN and dramatically increases the
number of valid target sites---typically at the cost of reduced specificity.  All four are
supported by swapping the two-character PAM check in the scanning loop for an IUPAC matcher.

\end{document}
